\documentclass[11pt]{article}
\usepackage[margin=1in]{geometry}
\usepackage{amsmath,amssymb,amsthm,mathtools,enumitem,hyperref}
\hypersetup{colorlinks=true,linkcolor=blue,urlcolor=blue,citecolor=blue}
\newtheorem{theorem}{Theorem}
\newtheorem{lemma}[theorem]{Lemma}
\newcommand{\eps}{\varepsilon}
\newcommand{\fall}[2]{(#1)_{\underline{#2}}}
\title{Weak virial positivity at every fixed order\\\large MathDB 351254}
\author{Principia Math In Progress}
\date{September 18, 2026}
\begin{document}
\maketitle
\begin{abstract}
For every fixed degree $r\ge1$, index $i\ge0$, and order $k\ge2$,
a uniform random simple $r$-regular bipartite graph on two labelled
parts of size $n$ satisfies $\Delta^k[-\log(j!m_j)]_{j=i}>0$ with
probability $1-O(n^{-2})$. This proves the full weak virial-positivity
conjecture. The argument combines the exact Wanless--Pernici matching
expansion and bounded contributor moments with an all-order formal
coefficient calculation, whose proof is included.
\end{abstract}

\section*{Task 1. Full solution}
\subsection*{Statement, model, and established inputs}
Let $m_j=m_j(G)$ be the number of $j$-matchings and let
$\Delta v(j)=v(j+1)-v(j)$. All parameters $r,i,k$ below are fixed
as $n\to\infty$. The random graph is uniform among simple
$r$-regular bipartite graphs with a fixed labelled bipartition of
sizes $n,n$, the standard model underlying the source's random
regular graph estimates. Write $u(j)=-\log(j!m_j)$.

\begin{theorem}\label{thm:main}
For every $r\ge1$, $i\ge0$, and $k\ge2$, there is a constant
$C_{r,i,k}$ such that, for all sufficiently large $n$,
\[
 \Pr\{\Delta^k u(i)\le0\}\le C_{r,i,k}n^{-2}.
\]
In particular, Conjecture I.2 of \cite{virial} holds without the
restrictions in that paper's Theorem I.1.
\end{theorem}

We use the following two established facts in the precise form
recorded in \cite[Section 3, equations (3.1)--(3.10)]{graph}.
Set
\begin{align}
 T_r(y)&=\frac{r\sqrt{1-4(r-1)y}-r+2}{2(1-r^2y)}
        =\sum_{s\ge0}t_s y^s,\label{eq:T}\\
 f(x)&=-\sum_{s\ge1}\frac{t_s}{s}(-x)^s,
 &M_j(n)&=[x^j]e^{nf(x)}.\label{eq:M}
\end{align}
Here $t_0=1$, $t_1=r$, and the square root has constant term one.
The rationalized form defines $r=1$ as well.

First, the exact Wanless--Pernici expansion \cite{wanless,pernici} is
\begin{equation}\label{eq:matching}
 m_j=\sum_{z=0}^j c_z(G)M_{j-z}(n),\qquad
 \sum_{z\ge0}c_z x^z
   =\exp\left(\sum_{s\ge4}\frac{\eta_s(G)}{2s}(-x)^s\right).
\end{equation}
In particular, $c_0=1$, $c_1=c_2=c_3=0$. Each $\eta_s$ is a
fixed linear combination, with coefficients depending on $r$, of
counts of finitely many fixed cyclic contributor subgraphs.
Second, for every fixed finite collection of these variables and
every fixed polynomial $P$,
\begin{equation}\label{eq:moments}
 \sup_{n\text{ sufficiently large}}\mathbb E|P(\eta_4,\ldots,\eta_J)|^2
 <\infty.
\end{equation}
This is the bounded-moment input of \cite[equation (3.10)]{graph};
its use is independent of that paper's finite-order coefficient
restrictions.

One can see the moment bound directly in the bipartite pairing
model. The expected number of embeddings of a fixed graph $H$ with
$v$ vertices and $e$ edges is $O(n^{v-e})$: there are $O(n^v)$
vertex choices, boundedly many stub choices, and any $e$ prescribed
pairs have probability $1/\fall{rn}{e}$. In a product of cyclic
contributor counts, classify the finitely many possible overlaps.
Every component of each union contains a cycle, so $e\ge v$.
All fixed moments are consequently bounded. Conditioning on
simplicity preserves these bounds, since its probability tends to
a positive constant for fixed $r$; this standard random-regular
graph fact is also among the inputs cited in \cite{graph}.
Polynomial expansion and Cauchy--Schwarz give \eqref{eq:moments}.

\subsection*{The coefficient lemma, at all orders}
The next lemma is the coefficient result established in our companion
MathDB 349137 solution. We reproduce its proof, including the known
logarithmic degree bound, to specify every algebraic step used here.

\begin{lemma}\label{lem:coeff}
Put $\eps=1/n$ and $A_j=j!M_j/(nr)^j$, with its polynomial extension
in $j$. For arbitrary formal constants $c_z$, $z\ge4$, define
\[
 F_j=A_j+\sum_{z\ge4}\frac{c_z}{r^z}\eps^z\fall{j}{z}A_{j-z}.
\]
Then, for all $h\ge1$,
\begin{equation}\label{eq:coeff}
 \deg_j[\eps^h]\log F_j\le h+1,\qquad
 [j^{h+1}\eps^h]\log F_j=\frac{r^{-h}-2}{h(h+1)}.
\end{equation}
All identities are coefficientwise formal identities, with the
$c_z$ independent of $j$ and $\eps$.
\end{lemma}
\begin{proof}
Write $f(x)=\sum_{s\ge1}f_sx^s$, where $f_1=r$.
Expansion of the exponential gives
\begin{equation}\label{eq:ah}
 [\eps^h]A_j=
 \sum_{\sum_{s\ge2}(s-1)m_s=h}
 \fall{j}{\sum s m_s}\prod_{s\ge2}
       \frac{(f_s/r^s)^{m_s}}{m_s!}.
\end{equation}
This is a polynomial of degree at most $2h$.

For the stronger logarithmic bound, view $A_j$ as the partition
function of disjoint subsets $B\subseteq[j]$, $|B|=s\ge2$, with
weight $s!f_sr^{-s}\eps^{s-1}$. Its logarithm is the sum of
connected clusters of intersecting subsets. Explicitly, the weight
of an ordered tuple $B_1,\ldots,B_m$ in the logarithm is its product
weight divided by $m!$, multiplied by
\[
 \sum_{H\text{ connected on }[m]}
 \prod_{\{a,b\}\in E(H)}(-\mathbf1_{B_a\cap B_b\ne\varnothing}).
\]
This identity follows by expanding the incompatibility product over
all graphs and decomposing each graph into connected components;
repeated subsets in the tuples are allowed. Every weight has
positive $\eps$-order. A connected term of order
$h=\sum_a(|B_a|-1)$ uses at most $h+1$ labels, by adding subsets
along a spanning tree of its intersection graph. Grouping terms by
their union size therefore expresses
\[
 L_j:=\log A_j=\sum_{h\ge1}\ell_h(j)\eps^h,
 \qquad \deg\ell_h\le h+1.
\]
This is the degree bound already proved in \cite[Theorem 3.1]{graph}.

Put $B(t)=\sum_{h\ge1}[j^{h+1}]\ell_h(j)\,t^{h+1}$.
Differentiating $e^{nf(x)}$ gives, after normalization,
\begin{equation}\label{eq:rec}
 j\eps=\sum_{s\ge1}\frac{sf_s}{r^s}\eps^s\fall{j}{s}
                  \exp(L_{j-s}-L_j).
\end{equation}
The identity first holds for nonnegative integer $j$. At each
$\eps$-order it is polynomial in $j$, so it holds formally in $j$.
Substitute $j=t/\eps$. Fixed shifts lower polynomial degree:
$\deg(\ell_h(j-s)-\ell_h(j))\le h$. Hence the shifted logarithm
belongs to $\mathbb Q(r)[[t,\eps]]$, and its $\eps^0$ coefficient
is $-sB'(t)$. Also
$\eps^s\fall{j}{s}=\prod_{a=0}^{s-1}(t-a\eps)$ has total
degree $s$. Thus taking the $\eps^0$ coefficient in \eqref{eq:rec}
is legitimate and yields
\[
 t=xf'(x),\qquad x=\frac tr e^{-B'(t)}.
\]
From \eqref{eq:T}--\eqref{eq:M},
\[
 xf'(x)=1-T_r(-x),\qquad
 (1-r^2y)T_r(y)^2+(r-2)T_r(y)-(r-1)=0.
\]
Putting $T_r(-x)=1-t$ therefore gives
\[
 x=\frac{t(r-t)}{r^2(1-t)^2},\qquad
 B'(t)=\log\frac{t}{rx}
       =2\log(1-t)-\log(1-t/r).
\]
The coefficient in \eqref{eq:coeff} follows for $A_j$.

Finally, series $\sum_h p_h(j)\eps^h$ with $\deg p_h\le h$
form a ring closed under formal exponentials with zero constant
term and logarithms with constant term one. Each $L_{j-z}-L_j$
is in this ring. Therefore
\[
 F_j/A_j=1+\sum_{z\ge4}\frac{c_z}{r^z}
      \eps^z\fall{j}{z}\exp(L_{j-z}-L_j)
\]
is in the ring as well, and $[\eps^h]\log(F_j/A_j)$ has degree
at most $h$. This proves both assertions for $F_j$.
\end{proof}

\subsection*{A finite polynomial with positive leading term}
Fix $J=i+k$ and set $B_j=j!m_j/(nr)^j$ for $0\le j\le J$.
Treating $c_4,\ldots,c_J$ as indeterminates, \eqref{eq:matching}
identifies $B_j$ with $F_j$. For integer $j$, \eqref{eq:ah} is
the exact polynomial $A_j$, and terms with $z>j$ vanish. Thus every
$B_j$ is a finite polynomial in $\eps$ and the $c_z$, with
constant term one.

Let $s_l=(-1)^{k-l}$ and define finite polynomials
\[
 P_+=\prod_{s_l=1}B_{i+l}^{\binom kl},\qquad
 P_-=\prod_{s_l=-1}B_{i+l}^{\binom kl},\qquad D=P_--P_+.
\]
The $k$th difference kills polynomials of degree less than $k$ and
sends $j^k$ to $k!$. Lemma \ref{lem:coeff} gives, in the formal
polynomial coefficient ring,
\begin{equation}\label{eq:logdifference}
 \log P_+-\log P_-=
 -a_{r,k}\eps^{k-1}+O(\eps^k),\qquad
 a_{r,k}=(k-2)!\bigl(2-r^{-(k-1)}\bigr)>0.
\end{equation}
Since $P_-(0)=1$, exponentiation yields the exact finite expansion
\begin{equation}\label{eq:D}
 D=a_{r,k}\eps^{k-1}+\sum_{h=k}^{H}R_h(c_4,\ldots,c_J)\eps^h
\end{equation}
for some finite $H$ and fixed polynomials $R_h$. The formal
calculation determines the initial coefficients of a finite
polynomial; it makes no convergence claim about a random logarithm.

Now evaluate the $c_z$ on the random graph. Each $R_h$ is a fixed
polynomial in finitely many $\eta_s$, by \eqref{eq:matching}, so
\eqref{eq:moments} bounds its second moment. Equation \eqref{eq:D}
and the triangle inequality in $L^2$ give
\[
 \bigl\|n^{k-1}D-a_{r,k}\bigr\|_2
 \le\sum_{h=k}^H n^{k-1-h}\|R_h\|_2=O(n^{-1}).
\]
Consequently,
\begin{equation}\label{eq:prob}
 \Pr\{D\le0\}\le
 \frac{\mathbb E|n^{k-1}D-a_{r,k}|^2}{a_{r,k}^2}=O(n^{-2}).
\end{equation}

A regular bipartite graph has a perfect matching by Hall's theorem:
for each set of left vertices, counting its incident edges gives
$r|S|\le r|N(S)|$. Thus $m_j>0$ for $0\le j\le n$. For $n\ge J$
all the $B_j,P_+,P_-$ used above are positive. Since $k\ge2$,
the term $j\log(nr)$ disappears under $\Delta^k$, and
\[
 \Delta^k u(i)=\log P_--\log P_+.
\]
Its sign equals that of $D$. The theorem follows from \eqref{eq:prob}.
\hfill$\square$

\subsection*{Attribution, verification, and scope}
The matching expansion and random-graph moment bounds are established
inputs, credited above. The leading coefficient is Pernici's formula;
the all-order proof is reproduced from our MathDB 349137 result.
The finite-product probability method follows Federbush. The new
conclusion here is its unrestricted application to every fixed
$r,i,k$, with the explicit failure bound. This does not establish
simultaneous positivity for indices growing with $n$.

The exact verifier checks the source coefficients, logarithmic
degree and difference cancellations, finite-product leading terms,
and matching counts obtained independently from graph enumeration.
All five pages of the current primary and all 18 pages of its
companion proof paper were read. Bounded checks found no target
journal final or earlier external full claim. Review is by the
authoring agent, without independent peer review or proof-assistant
certification.

\section*{Tasks 2--5. Proposed follow-up directions}
\begin{enumerate}[label=\textbf{Task \arabic*.},start=2,leftmargin=*]
\item Determine explicit constants in the failure bound in terms of
$r,i,k$, using contributor-count moments and the finite polynomial
in \eqref{eq:D}.
\item Use higher bounded moments to sharpen the fixed-parameter
failure probability, and identify the first nontrivial fluctuation
of the scaled polynomial $n^{k-1}D$.
\item Obtain estimates uniform in increasing $i$ and $k$, to approach
the separate strong virial-positivity conjecture.
\item Apply the same all-order coefficient calculation to weak graph
positivity and to the factorial-modified matching quantity discussed
at the end of \cite{virial}.
\end{enumerate}
\begin{thebibliography}{9}
\bibitem{virial} P. Federbush, \emph{Random Regular Bipartite Graphs
Satisfy Weak Virial Positivity, for a Large Range of the Parameters},
\href{https://arxiv.org/abs/2107.05110v2}{arXiv:2107.05110v2}, 2022.
\bibitem{graph} P. Federbush, \emph{A PROOF of Weak Graph Positivity,
for a Large Range of the Parameters},
\href{https://arxiv.org/abs/2012.10927v8}{arXiv:2012.10927v8}, 2022.
\bibitem{wanless} I. M. Wanless, \emph{Counting matchings and tree-like
walks in regular graphs}, Combin. Probab. Comput. \textbf{19} (2010), 463.
\bibitem{pernici} M. Pernici, \emph{$1/n$ expansion for the number of
matchings on regular graphs and monomer-dimer entropy},
J. Stat. Phys. \textbf{168} (2017), 666.
\end{thebibliography}
\end{document}
